\clearpage\thispagestyle{empty}
\pdfbookmark[0]{Quellen zu Komet (Oktober 1811)}{comet1811oct-sources}
\begin{center}\textbf{Quellen zu Komet (Oktober 1811)}\end{center}
\noindent Diese KI-Übertragung ins heutige Deutsch folgt dem redaktionellen Bericht und dem Auszug aus dem Brief von Carl Friedrich Gauß vom 14. Oktober 1811 in \href{https://gdz.sub.uni-goettingen.de/id/PPN236007467}{\textit{Werke}, Band~VI} (Göttingen, 1874), S.~338--340. Der Erstdruck erschien als Beitrag XL, \textit{Über den großen Cometen von 1811}, in der \textit{Monatlichen Correspondenz}, Band XXIV, Oktober 1811, S.~\href{https://api.digitale-sammlungen.de/iiif/image/v2/bsb10538613_00418/full/1600,/0/default.jpg}{406}, \href{https://api.digitale-sammlungen.de/iiif/image/v2/bsb10538613_00419/full/1600,/0/default.jpg}{407}, \href{https://api.digitale-sammlungen.de/iiif/image/v2/bsb10538613_00420/full/1600,/0/default.jpg}{408}, \href{https://api.digitale-sammlungen.de/iiif/image/v2/bsb10538613_00421/full/1600,/0/default.jpg}{409}, \href{https://api.digitale-sammlungen.de/iiif/image/v2/bsb10538613_00422/full/1600,/0/default.jpg}{410} und \href{https://api.digitale-sammlungen.de/iiif/image/v2/bsb10538613_00423/full/1600,/0/default.jpg}{411}. Er bezeichnet sich als Fortsetzung des Beitrags, der im Septemberheft auf S.~289 beginnt. Das digitalisierte Exemplar stammt aus der Bayerischen Staatsbibliothek.
\par\smallskip\noindent In der Zeitschrift steht das Briefdatum in einer Fußnote auf S.~409. Die Werkausgabe nimmt das Datum in den Text auf und endet mit Gauß’ Bemerkung, Olbers werde ausführlicher schreiben. In der Zeitschrift folgt darauf die Fortsetzung des redaktionellen Berichts mit dem Beginn der Vergleiche mit von Zachs Beobachtungen in Marseille. Die Zeitschrift setzt in der Deklinationsspalte am 10. September einen Doppelstrich zu $-27$, verwendet beim Periheldurchgang andere Zeitzeichen, gibt die Winkeleinheiten der Bahnneigung an und quadriert den Kosinus in allen drei Koordinatenformeln. Zeichen, Einheiten und Nenner folgen in der Übertragung der Werkausgabe. KI-Übertragung und KI-Satz.
\clearpage
